
The baseline controller of \cref{sec:proposed} achieves fault
tolerance through a single structural choice: setting the altitude
thrust feed-forward equal to the drone's own measured rope tension,
$T_i^{\mathrm{ff}} = T_i$. We establish that this identity renders
the closed loop \emph{hybrid practical input-to-state stable}
(hybrid practical-ISS) under the cable-severance fault model of
\cref{sec:reduction-and-reference}. When a peer cable severs,
$T_i$ rises immediately and the feed-forward converts that step
into a proportional thrust increment at the plant's native rate,
without any detection, annunciation, or reconfiguration interval.
The argument proceeds through the analytical hypotheses, six
lemmas, the main theorem, and a claims-to-evidence map.

\paragraph{Status of the lemma stack.}
The six lemmas and the main theorem certify the architectural
claim of \cref{sec:Introduction}---that within the analyzed
baseline architecture and its declared admissibility domain, a
single decentralized feed-forward identity supplies the
bounded-jump mechanism required for hybrid practical-ISS---rather
than constituting a stand-alone theoretical contribution. The
hybrid-systems framework is borrowed
from~\cite{GoebelSanfeliceTeel2012}; the practical-ISS
composition tools are standard. The value added is showing that
the framework's hypotheses hold for the proposed architecture on
the declared admissibility domain $\Omega_\tau^{\mathrm{dwell}}$,
with closed-form expressions for $\rho$ and $\lambda$, and a
calibrated class-$\mathcal{K}$ bound for the per-fault jump
$\chi(\Delta_f)$.

% ----------------------------------------------------------
\subsection{Setup, Notation, and Analytical Hypotheses}
\label{sec:stability-setup}
% ----------------------------------------------------------

The analysis operates on the reduced-order model of
\cref{sec:reduction-and-reference}. Let
$e_{p,z,i}(t) = p_{\mathrm{slot},i,z}(t) - p_{i,z}(t)$ be
drone $i$'s altitude slot-tracking error and
$\vect{e}_p(t) \in \R^{3N_s}$ the stacked position error over
$\mathcal{S}(t)$ with $N_s = |\mathcal{S}(t)|$. Fault events
and the dwell condition follow \cref{sec:fault-and-controller};
consecutive faults are separated by at least one
payload-pendulum period,
\begin{equation}
  \tau_{d,k} \;\ge\; \tau_{\mathrm{pend}}
    \;\triangleq\; 2\pi\sqrt{L/g}.
  \label{eq:dwell-condition}
\end{equation}
At $L = \SI{1.25}{\metre}$, $\tau_{\mathrm{pend}} =
\SI{2.24}{\second}$ and the pendulum natural frequency is
$\omega_p \triangleq \sqrt{g/L} = \SI{2.80}{\radian\per\second}$.

The analysis rests on three hypotheses, defended empirically in
\cref{tab:stability-evidence,tab:domain-audit}.
\textbf{H1}~(slack-excursion budget) requires every trajectory
to remain in $\Omega_\tau^{\mathrm{dwell}}$: maximum slack-run
duration $\tau_{\mathrm{slack,max}} \le \SI{40}{\milli\second}$
and duty cycle $\eta_{\max} \le 2.5\%$.
\textbf{H2}~(spaced-fault dwell) requires
$\tau_{d,k} \ge \tau_{\mathrm{pend}}$~\eqref{eq:dwell-condition}.
\textbf{H3}~(Doctrine~D1, single-active-regime QP) requires the
per-tick QP's active set $\mathcal{A}(t)$ to be locally constant
inside each inter-fault window; transitions are absorbed into the
H1 $\eta_{\max}$ budget and empirically hold on $\ge 99.9\%$ of
mission time (worst transition fraction $0.10\%$,
\cref{tab:domain-audit}). Claims outside H1--H3---rapid
multi-fault, long slack excursions, QP-fallback---are out of
scope; the analysis produces no bounds for trajectories
violating them.

\paragraph{Partial circularity and its mitigation.}
The empirical defence of H1 and H3 uses the same pre-declared
missions as the theorem's consequences, so the validation is
partially \emph{post hoc}. The theorem therefore provides a
\emph{conditional} guarantee: it holds whenever the trajectory
satisfies H1--H3, and the margins of \cref{tab:domain-audit}
document how far inside the domain the demonstrated missions
sit. Two measures mitigate the circularity. First, the gate
bounds ($\tau_{\mathrm{slack,max}} = \SI{40}{\milli\second}$,
$\eta_{\max} = 2.5\%$) are set by the physics of bead-chain
re-tension, not by data-fitting: $\tau_{\mathrm{slack,max}}$ is
one quarter of the Tikhonov frozen-window $\tau_{\mathrm{pend}}
/(2\pi)$ and $\eta_{\max}$ is a Rayleigh-turbulence budget.
Second, the taut-cable reduction (\cref{thm:reduction}, C1) is
independently validated against the bead-chain truth model in
\cref{sec:VI-reduction}, providing an \emph{a~priori} check
that the reduced model tracks the full plant within
$O(\delta + \eta_{\max})$ on $\Omega_\tau^{\mathrm{dwell}}$;
the constants $C_1, C_2$ are bounded by the bead-chain
parameters $(m_{\mathrm{bead}}, c_s, k_s, N_b)$ and the wind
bound $W_{\max}$ via the Tikhonov singular-perturbation argument
of \cref{sec:reduction-and-reference}, and the stability
analysis absorbs the residual into $\gamma_\eta$ of
\cref{thm:hybrid-stability} without appeal to the simulated
missions.
Design-time certification of H1 and H3 for unseen trajectory
families---which would require trajectory-planning machinery
beyond the present scope---is identified as future work.

% ----------------------------------------------------------
\subsection{Lemmas Underpinning the Hybrid Practical-ISS Theorem}
\label{sec:five-lemmas}
% ----------------------------------------------------------

We state six lemmas that support the main theorem.
\Cref{lem:tension-cancel} establishes algebraic tension
cancellation on the altitude channel;
\cref{lem:altitude-iss} gives exponential ISS for the altitude
channel through the closed-loop Lyapunov matrix $P_v$;
\cref{lem:horizontal-iss} gives the parallel result for the
horizontal channels through $P_{xy}$, providing the
$\alpha_{xy}$ rate that drives the dwell-cycle contraction
$\rho$;
\cref{lem:antiswing-invariant} shows the anti-swing damping is
invariant to surviving cardinality;
\cref{lem:jump-continuity} gives a bounded Lyapunov jump at
each fault in terms of the quasi-static redistribution
magnitude $\Delta_f$, together with the equal-sharing
specialization used for canonical calibration;
\cref{lem:dwell-decay} composes jump and decay into the
dwell-cycle contraction $\rho \in (0, 1)$.

For fault jumps, the tension coordinates are compared on the
post-fault survivor set by restricting the pre-fault vector to
the remaining cables; this convention is only needed in
\cref{lem:jump-continuity}.

% ---- Lemma 2 (unchanged) ----------------------------------
\begin{lemma}[Tension cancellation]
\label{lem:tension-cancel}
Assume H1 and the cable-tilt domain condition
$|\theta_i(t)| \le \theta_{c,\max}$ for known
$\theta_{c,\max} \in (0, \pi/2)$. Then
$T_i^{\mathrm{ff}} = T_i$ cancels the dominant projection term
$T_i\cos\theta_i$ from the drone altitude acceleration,
leaving a residual factor $(1 - \cos\theta_i)$ bounded by
$1 - \cos\theta_{c,\max}$. At $\theta_{c,\max} =
\SI{0.75}{\radian}$ this is $0.27$, and the residual enters
the altitude-error equation as a bounded matched disturbance
handled by \cref{lem:altitude-iss}. The bound holds uniformly
in $N_s \ge 2$.
\end{lemma}

\begin{proof}
Substituting $f_i = m_i(g + a_{d,z}^\star) + T_i^{\mathrm{ff}}$
into $m_i \ddot{p}_{i,z} = f_i - m_i g - T_i\cos\theta_i +
F_{w,i,z}$ and absorbing the wind term yields
\begin{equation}
  m_i \ddot{p}_{i,z}
    = m_i a_{d,z}^\star + T_i(1 - \cos\theta_i).
  \label{eq:alt-dynamics-residual}
\end{equation}
Since $|\theta_i| \le \theta_{c,\max}$ uniformly (enforced by
the QP tilt ceiling, \cref{sec:qp-and-ff}), the residual
$T_i(1-\cos\theta_i) \le T_i(1 - \cos\theta_{c,\max})$
regardless of $N_s$, and is absorbed by $\gamma_q$ of
\cref{lem:altitude-iss}.
\end{proof}

% ---- Lemma 3 (altitude ISS — expanded proof) --------------
\begin{lemma}[Altitude ISS via Lyapunov matrix]
\label{lem:altitude-iss}
Under hypotheses H1--H3, the altitude tracking error
$e_{p,z,i}$ satisfies exponential input-to-state stability
with explicit gain functions: there exist $\lambda_z > 0$
and $\gamma_w, \gamma_q \in \mathcal{K}$ such that, along
trajectories between fault events,
\begin{equation}
  \dot V_z \;\le\; -\lambda_z\,V_z
    \;+\; \gamma_w(\|w\|_\infty)
    \;+\; \gamma_q(\|q - q^\star\|_\infty),
  \label{eq:lyapunov-decay-alt}
\end{equation}
where $V_z(e,\dot e) \triangleq \tfrac{1}{2}
[e,\dot e]^\top P_v\,[e,\dot e]$ and $P_v \succ 0$ is the
unique positive-definite solution of
$A_m^\top P_v + P_v A_m = -I$ for the closed-loop matrix
$A_m$. The ISS decay rate is
$\lambda_z = \alpha_z \triangleq |\mathrm{Re}(s_1)| =
\SI{5.37}{\second^{-1}}$, where $s_1$ is the dominant pole
of $A_m$.
\end{lemma}

\begin{proof}
After tension cancellation (\cref{lem:tension-cancel}), the
altitude slot-tracking error $e_{p,z,i} \triangleq
p_{\mathrm{slot},i,z} - p_{i,z}$ satisfies
\begin{equation}
  \ddot{e}_{p,z,i}
    + K_d^z \dot{e}_{p,z,i}
    + K_p^z e_{p,z,i}
    = \ddot{p}_{\mathrm{slot},z} + r_i(t),
  \label{eq:alt-error-ode}
\end{equation}
with residual $r_i(t) = O\!\left(\theta_{c,\max}^2
T_i/m_i\right) + O(\delta)$ bounded uniformly on
$\Omega_\tau^{\mathrm{dwell}}$ under H1. The closed-loop
matrix
\begin{equation}
  A_m = \begin{bmatrix}0 & 1\\ -K_p^z & -K_d^z\end{bmatrix}
      = \begin{bmatrix}0 & 1\\ -100 & -24\end{bmatrix}
\end{equation}
has characteristic polynomial
$s^2 + 24s + 100 = 0$, yielding eigenvalues
\begin{equation}
  s_{1,2}
    = \frac{-24 \pm \sqrt{576 - 400}}{2}
    = -12 \pm \sqrt{44},
\end{equation}
so $s_1 = -12 + \sqrt{44} \approx -5.37\,\si{\second^{-1}}$
(dominant) and $s_2 \approx -18.63\,\si{\second^{-1}}$.
Both eigenvalues have strictly negative real parts, so $A_m$
is Hurwitz and the unique solution to
$A_m^\top P_v + P_v A_m = -I$ is $P_v \succ 0$:
\begin{equation}
  P_v = \begin{bmatrix}
    p_{11} & p_{12}\\ p_{12} & p_{22}
  \end{bmatrix}.
  \label{eq:Pv-generic}
\end{equation}
Expanding $A_m^\top P_v + P_v A_m = -I$ at $K_p^z = 100$,
$K_d^z = 24$ gives the three scalar equations
$-2K_p^z p_{12} = -1$, $p_{11} - K_p^z p_{22} - K_d^z p_{12}
= 0$, and $2 p_{12} - 2 K_d^z p_{22} = -1$, whence
\begin{equation}
  p_{12} = \frac{1}{2K_p^z} = \frac{1}{200} = 0.005,\quad
  p_{22} = \frac{1 + 2 p_{12}}{2K_d^z}
         = \frac{1.01}{48} \approx 0.0210,\quad
  p_{11} = K_p^z p_{22} + K_d^z p_{12} \approx 2.224,
\end{equation}
yielding the numerical matrix
\begin{equation}
  P_v = \begin{bmatrix} 2.224 & 0.005 \\ 0.005 & 0.0210
        \end{bmatrix},
  \qquad \lambda_{\min}(P_v) \approx 0.021,\quad
         \lambda_{\max}(P_v) \approx 2.224.
  \label{eq:Pv-numeric}
\end{equation}
Setting $\zeta \triangleq [e_{p,z,i},\,\dot{e}_{p,z,i}]^\top$
and $V_z = \tfrac{1}{2}\zeta^\top P_v \zeta$, differentiating
along trajectories of \eqref{eq:alt-error-ode} and applying
Young's inequality to the cross-terms from
$\ddot{p}_{\mathrm{slot},z} + r_i(t)$ yields
\eqref{eq:lyapunov-decay-alt} with decay rate $\lambda_z =
\alpha_z = |\mathrm{Re}(s_1)| = \SI{5.37}{\second^{-1}}$; the
ISS gains $\gamma_w, \gamma_q$ are proportional to
$\lambda_{\max}(P_v)/\lambda_{\min}(P_v) \approx 106$ and to
the magnitudes of $w$ and $r_i$ from \cref{lem:tension-cancel}
and $\delta$ from \cref{thm:reduction}. Concretely, the
Young's-inequality bound $2\zeta^\top P_v r_i \le
\varepsilon\,\zeta^\top P_v\zeta +
\kappa(P_v)\,\varepsilon^{-1}\|r_i\|^2$ (with
$\kappa(P_v) = \lambda_{\max}(P_v)/\lambda_{\min}(P_v)\approx
106$ and $\varepsilon = \alpha_z$) absorbs $r_i$ into a
$\mathcal{K}$-gain at the cost of the condition number.
The decay rate
$\lambda_z$ is independent of $N_s$ because $A_m$ depends only
on the per-drone outer-loop gains $(K_p^z, K_d^z)$.
\end{proof}

% ---- Lemma 3b (horizontal ISS — NEW) ----------------------
\begin{lemma}[Horizontal ISS via Lyapunov matrix]
\label{lem:horizontal-iss}
Under hypotheses H1--H3, the horizontal slot-tracking error
$e_{p,xy,i} \in \R^2$ satisfies exponential ISS: there exist
$\lambda_{xy} > 0$ and $\gamma_w^{xy}, \gamma_q^{xy} \in
\mathcal{K}$ such that, between fault events,
\begin{equation}
  \dot V_{xy} \;\le\; -\lambda_{xy}\,V_{xy}
    \;+\; \gamma_w^{xy}(\|w\|_\infty)
    \;+\; \gamma_q^{xy}(\|q - q^\star\|_\infty),
  \label{eq:lyapunov-decay-xy}
\end{equation}
where $V_{xy}(e,\dot e) \triangleq \tfrac{1}{2}
[e,\dot e]^\top P_{xy}\,[e,\dot e]$ and $P_{xy} \succ 0$ is
the unique positive-definite solution of
$A_{m,xy}^\top P_{xy} + P_{xy} A_{m,xy} = -I$ for the
horizontal closed-loop matrix $A_{m,xy}$.
The ISS decay rate satisfies
$\lambda_{xy} = \alpha_{xy} \triangleq
|\mathrm{Re}(s_1^{xy})| = \SI{2.38}{\second^{-1}}$,
where $s_1^{xy}$ is the dominant pole of $A_{m,xy}$.
Since $\alpha_{xy} < \alpha_z$, the horizontal rate is
$\alpha_{\min} = \min(\alpha_z,\alpha_{xy}) = \alpha_{xy}$,
which drives the dwell-cycle contraction of
\cref{lem:dwell-decay}.
\end{lemma}

\begin{proof}
The horizontal slot-tracking error in each of the $x$- and
$y$-channels satisfies an identical scalar second-order ODE:
\begin{equation}
  \ddot{e}_{p,xy,i}
    + K_d^{xy} \dot{e}_{p,xy,i}
    + K_p^{xy} e_{p,xy,i}
    = \ddot{p}_{\mathrm{slot},xy} + r_{xy,i}(t),
\end{equation}
where $r_{xy,i}(t) = O(\theta_{c,\max} T_i/m_i) + O(\delta)$
is bounded by H1 and \cref{lem:tension-cancel}. No tension
feed-forward enters the horizontal channels, so
\cref{lem:tension-cancel} does not directly apply here; the
horizontal residual is the tilt-projection component
$T_i \sin\theta_i$ redistributed through the attitude loop.
The horizontal closed-loop matrix is
\begin{equation}
  A_{m,xy} = \begin{bmatrix}0 & 1 \\ -K_p^{xy} & -K_d^{xy}
              \end{bmatrix}
           = \begin{bmatrix}0 & 1 \\ -30 & -15
              \end{bmatrix},
\end{equation}
with characteristic polynomial $s^2 + 15s + 30 = 0$ and
eigenvalues
\begin{equation}
  s_{1,2}^{xy}
    = \frac{-15 \pm \sqrt{225 - 120}}{2}
    = \frac{-15 \pm \sqrt{105}}{2},
\end{equation}
yielding $s_1^{xy} = (-15 + \sqrt{105})/2 \approx
-2.38\,\si{\second^{-1}}$ (dominant) and
$s_2^{xy} \approx -12.625\,\si{\second^{-1}}$.
The unique solution to
$A_{m,xy}^\top P_{xy} + P_{xy} A_{m,xy} = -I$
is obtained by solving the three scalar equations
$-2K_p^{xy} p_{12}^{xy} = -1$,
$2p_{12}^{xy} - 2K_d^{xy} p_{22}^{xy} = -1$, and
$p_{11}^{xy} = K_p^{xy} p_{22}^{xy} + K_d^{xy} p_{12}^{xy}$:
\begin{equation}
  p_{12}^{xy} = \frac{1}{2K_p^{xy}} = \frac{1}{60} \approx 0.0167,\quad
  p_{22}^{xy} = \frac{1 + 2p_{12}^{xy}}{2K_d^{xy}}
              = \frac{31}{900}
              \approx 0.0344,\quad
  p_{11}^{xy} = K_p^{xy} p_{22}^{xy} + K_d^{xy} p_{12}^{xy}
              = \frac{77}{60} \approx 1.283,
\end{equation}
giving
\begin{equation}
  P_{xy} = \begin{bmatrix}
    1.283 & 0.0167 \\
    0.0167 & 0.0344
  \end{bmatrix},
  \qquad
  \lambda_{\min}(P_{xy}) \approx 0.034,\quad
  \lambda_{\max}(P_{xy}) \approx 1.283.
  \label{eq:Pxy-numeric}
\end{equation}
The Lyapunov argument of \cref{lem:altitude-iss} applies
verbatim to $V_{xy} = \tfrac{1}{2}\zeta_{xy}^\top P_{xy}
\zeta_{xy}$ with decay rate $\lambda_{xy} = \alpha_{xy} =
|\mathrm{Re}(s_1^{xy})| = \SI{2.38}{\second^{-1}}$.
Since $\alpha_{xy} = 2.38 < \alpha_z = 5.37$, we have
$\alpha_{\min} = \alpha_{xy}$, confirming the horizontal
channel as the bottleneck for the dwell-cycle contraction.
The ISS property is independent of $N_s$ for the same reason
as in \cref{lem:altitude-iss}.
\end{proof}

% ---- Lemma 4 (unchanged) ----------------------------------
\begin{lemma}[Pendulum damping invariance]
\label{lem:antiswing-invariant}
The active anti-swing damping coefficient
$B_{\mathrm{swing}}$ is independent of the number of
surviving drones $N_s$, ensuring that the pendulum damping
ratio $\zeta > 1$ is preserved across any cable-severance
fault transition.
\end{lemma}

\begin{proof}
The restoring force is
$F_{\mathrm{restore}} = N_s T_i^{\mathrm{ss}}
k_{\mathrm{swing}} v_{L,\perp}/L$
with $T_i^{\mathrm{ss}} = m_L g / N_s$, so
$N_s T_i^{\mathrm{ss}} = m_L g$ cancels the $N_s$ dependence
and leaves $B_{\mathrm{swing}} = m_L g\,k_{\mathrm{swing}}/L
= \SI{62.8}{\newton\second\per\metre}$.
With $\omega_p = \SI{2.80}{\radian\per\second}$, the damping
ratio $\zeta = B_{\mathrm{swing}}/(2 m_L \omega_p) = 1.12 >
1$; the pendulum is overdamped for any $N_s \ge 2$.
\end{proof}

% ---- Lemma 5 (jump — expanded proof, explicit chi) --------
\begin{lemma}[Bounded fault jump with explicit $\chi$]
\label{lem:jump-continuity}
At each fault event $t_f$, the Lyapunov function undergoes a
bounded jump,
\begin{equation}
  V(\vect\xi(t_f^+))
    \;\le\;
    V(\vect\xi(t_f^-))
    \;+\; \chi(\Delta_f),
  \qquad
  \chi \in \mathcal{K},
  \label{eq:jump-bound}
\end{equation}
where $\vect\xi = (\vect e_p, \vect e_v, \vect\varepsilon_T)$,
the pre-fault tension coordinates are restricted to the
post-fault survivor set through the map
$\Pi_{\mathcal{S}^+\leftarrow\mathcal{S}^-}$, and
$\Delta_f = \|\vect T^{\mathrm{qs}}(t_f^+) -
\Pi_{\mathcal{S}^+\leftarrow\mathcal{S}^-}
\vect T^{\mathrm{qs}}(t_f^-)\|$ is the quasi-static tension
redistribution over the surviving cables induced by the
severance. On the admissible domain of H1--H3, the jump
function admits the bound
\begin{equation}
  \chi(\Delta_f)
    \;\le\;
    \bar c_f\bigl(\Delta_f + \Delta_f^2\bigr),
  \label{eq:chi-explicit}
\end{equation}
for some $\bar c_f > 0$ depending only on the admissible
domain and the Lyapunov weighting of the tension channel.
Under symmetric equal sharing at hover, the linear term
vanishes and \eqref{eq:chi-explicit} reduces to
$\chi(\Delta_f) \le \kappa_V\Delta_f^2/(2N_s)$.
\end{lemma}

\begin{proof}
The Kelvin--Voigt cable model keeps the kinematic state
$(\vect p_i, \vect v_i, p_L, v_L)$ continuous across $t_f$
since no impulsive force is applied; on the severed rope,
tension drops to zero within $\tau_{\mathrm{rope}} \approx
\SI{6.3}{\milli\second}$ (\cref{eq:tau-rope}).
Decompose
\begin{equation}
  V = V_{\mathrm{track}}(\vect e_p, \vect e_v)
    + \kappa_V\,V_{\mathrm{tension}}(\vect\varepsilon_T).
\end{equation}

\noindent\textit{Continuity of $V_{\mathrm{track}}$.}
Since $(\vect e_p, \vect e_v)$ is continuous at $t_f$
(the kinematic state is continuous and the slot reference
is pre-loaded), the survivor-restricted tracking term does
not increase at the fault instant, so
$V_{\mathrm{track}}(t_f^+) \le V_{\mathrm{track}}(t_f^-)$.

\noindent\textit{Jump in $V_{\mathrm{tension}}$.}
Let $\Pi_{\mathcal{S}^+\leftarrow\mathcal{S}^-}$ denote the
restriction map from the pre-fault survivor set to the
post-fault survivor set. The Kelvin--Voigt tension
$T_i = k_s(\ell_i - L)^+$ is a continuous functional of the
kinematic state $(\vect p_i, \vect v_i, \vect p_L, \vect v_L)$;
since no impulsive force is applied at $t_f$, the kinematic
state is continuous across the fault, so the survivor tensions
are exactly continuous at the fault instant; concretely,
\begin{equation}
  \vect\varepsilon_T(t_f^+)
    = \Pi_{\mathcal{S}^+\leftarrow\mathcal{S}^-}
      \vect\varepsilon_T(t_f^-)
      - \vect b_f,
  \qquad
  \vect b_f \triangleq
    \vect T^{\mathrm{qs}}(t_f^+)
    - \Pi_{\mathcal{S}^+\leftarrow\mathcal{S}^-}
      \vect T^{\mathrm{qs}}(t_f^-),
\end{equation}
with $\|\vect b_f\| = \Delta_f$. Since restriction cannot
increase the norm and the admissible domain of H1--H3 is
compact, $\|\vect\varepsilon_T(t)\|$ is uniformly bounded on
that domain. Since $V_{\mathrm{tension}}$ is quadratic,
expansion about
$\Pi_{\mathcal{S}^+\leftarrow\mathcal{S}^-}\vect\varepsilon_T(t_f^-)$
therefore gives finite constants $c_{f,1}, c_{f,2} > 0$ such
that
\begin{equation}
  V_{\mathrm{tension}}(t_f^+) - V_{\mathrm{tension}}(t_f^-)
    \;\le\;
    c_{f,1}\,\Delta_f + c_{f,2}\,\Delta_f^2.
\end{equation}
Absorbing $\kappa_V c_{f,1}$ and $\kappa_V c_{f,2}$ into
$\bar c_f$ gives \eqref{eq:chi-explicit}, so
\eqref{eq:jump-bound} follows with
$\chi(\Delta_f) = \bar c_f(\Delta_f + \Delta_f^2) \in
\mathcal{K}$.

\noindent\textit{Equal-sharing hover calibration.}
If the fault occurs from the canonical symmetric hover
operating point, then $\vect\varepsilon_T(t_f^-) = 0$ and
the reset is purely the quasi-static change in load sharing.
In that special case the linear term vanishes and the direct
survivor-count algebra reduces \eqref{eq:chi-explicit} to
$\chi(\Delta_f) \le \kappa_V\Delta_f^2/(2N_s)$.
\end{proof}

% ---- Lemma 6 (dwell-decay, expanded remark) ---------------
\begin{lemma}[Dwell-cycle contraction]
\label{lem:dwell-decay}
Let $t_k^\star$ and $t_{k+1}^\star$ be consecutive fault
events separated by dwell-time $\tau_{d,k} \ge
\tau_{\mathrm{pend}}$. Under the bounded fault jump of
\cref{lem:jump-continuity} and the ISS decays of
\cref{lem:altitude-iss,lem:horizontal-iss}, the Lyapunov
function satisfies a per-fault-cycle contraction
\begin{align}
  V(\vect\xi(t_{k+1}^{\star,-}))
    &\;\le\;
    \rho\,V(\vect\xi(t_k^{\star,-})) \;+\; c,
  \label{eq:dwell-contraction}\\
  \rho
    &\;\triangleq\;
    \exp\!\left(-\alpha_{\min}\,\tau_{\mathrm{pend}}\right),
  \quad \alpha_{\min} = \alpha_{xy} = \SI{2.38}{\second^{-1}},
  \label{eq:rho-def}
\end{align}
where
\begin{equation}
  c \;=\; \rho\,\chi(\Delta_f) \;+\;
          \frac{c_0}{\alpha_{\min}},
  \quad
  c_0 = \sup_{t}\!
    \bigl[\gamma_w(\|w(t)\|_\infty)
          + \gamma_q(\|q - q^\star\|_\infty)
          + \gamma_\eta(\eta_{\max})\bigr].
  \label{eq:c-def}
\end{equation}
At canonical parameters, $\rho \approx
\exp(-2.38 \times 2.24) \approx 0.005$.
\end{lemma}

\begin{proof}
Applying \cref{lem:jump-continuity} at $t_k^\star$:
\begin{equation}
  V(t_k^{\star,+}) \;\le\; V(t_k^{\star,-}) + \chi(\Delta_f).
  \label{eq:jump-step}
\end{equation}
Between $t_k^{\star,+}$ and $t_{k+1}^{\star,-}$, the flow
satisfies $\dot V \le -\alpha_{\min} V + c_0$ (combining
\cref{lem:altitude-iss,lem:horizontal-iss} through the
composite Lyapunov function $V = V_z + V_{xy} + \kappa_V
V_{\mathrm{tension}}$, using $\alpha_{\min} =
\min(\alpha_z, \alpha_{xy}) = \alpha_{xy}$). Integrating
over an interval of length $\tau \ge \tau_{\mathrm{pend}}$:
\begin{equation}
  V(t_{k+1}^{\star,-})
    \;\le\;
    e^{-\alpha_{\min}\tau}\bigl(V(t_k^{\star,-})
      + \chi(\Delta_f)\bigr)
    \;+\; \frac{c_0}{\alpha_{\min}}(1 - e^{-\alpha_{\min}\tau}).
\end{equation}
Setting $\rho = e^{-\alpha_{\min}\tau_{\mathrm{pend}}}$ and
upper-bounding $\tau$ by $\tau_{\mathrm{pend}}$ (the minimum
dwell, giving the worst-case exponential decay) yields
\eqref{eq:dwell-contraction}--\eqref{eq:c-def}.
Since $\alpha_{\min} = \alpha_{xy} > 0$ and
$\tau_{\mathrm{pend}} > 0$, we have $\rho \in (0,1)$.
\end{proof}

\begin{remark}[Analytical $\rho$ vs.\ empirical $\hat\rho$:
               the role of the additive constant $c$]
\label{rem:rho-gap}
\Cref{lem:dwell-decay} gives $\rho_{\mathrm{ana}} =
\exp(-\alpha_{xy}\tau_{\mathrm{pend}}) \approx 0.005$, yet
the Lyapunov-proxy measurements of \cref{sec:VI-D} yield
$\hat\rho_{\mathrm{V4}} = 0.20$ and
$\hat\rho_{\mathrm{V5}} = 0.045$.
The gap does not indicate a failure of the exponential decay
rate; it arises from the additive constant $c > 0$ in
\eqref{eq:dwell-contraction}. The empirical ratio is
\begin{equation}
  \hat\rho_k
    \;\triangleq\;
    \frac{V(t_{k+1}^{\star,-})}{V(t_k^{\star,-})}
    \;\le\;
    \rho_{\mathrm{ana}}
    \;+\; \frac{c}{V(t_k^{\star,-})},
  \label{eq:rho-hat-decomp}
\end{equation}
so $\hat\rho$ equals the true contraction rate $\rho_{\mathrm{ana}}$
only when $c \ll V(t_k^{\star,-})$. The inequality
\eqref{eq:rho-hat-decomp} is conservative: the additive constant
$c$ acts as a floor that is always non-negative, so the empirical
ratio $\hat\rho_k$ upper-bounds the true $\rho_{\mathrm{ana}}$
plus the disturbance contribution; it is not expected to be
tight except in the zero-disturbance limit. Under the Dryden
$\SI{4}{\metre\per\second}$ wind of V4--V5, the disturbance
floor $c_0/\alpha_{\min}$ constitutes a non-negligible fraction
of the pre-fault Lyapunov value: from \eqref{eq:c-def},
$c_0 \approx \gamma_w(W_{\max}) + \gamma_\eta(\eta_{\max})$
is set by the wind and reduction residuals, which are $O(1)$
in the Lyapunov proxy units. The ratio $c/V(t_k^{\star,-})
\approx 0.20 - 0.005 = 0.195$ for V4 and $\approx 0.040$
for V5 (where the longer dwell allows $V(t_2^{\star,-})$ to
reach a lower floor). The improvement from V4 to V5
($\hat\rho = 0.20 \to 0.045$) confirms \eqref{eq:rho-hat-decomp}:
with ten seconds of dwell instead of five, $V(t_2^{\star,-})$
has decayed further toward the disturbance floor, reducing the
$c/V$ ratio. In disturbance-free conditions ($w \equiv 0$,
$\delta \to 0$, $\eta_{\max} \to 0$), $c \to \rho\,\chi(\Delta_f)$
and $\hat\rho \to \rho_{\mathrm{ana}}$ asymptotically.
\end{remark}

% ----------------------------------------------------------
\subsection{Main Theorem, Robustness, and Multi-Fault Feasibility}
\label{sec:main-theorem-and-robustness}
% ----------------------------------------------------------

We now compose the six lemmas into the hybrid practical-ISS
theorem and state corollaries on robustness, the steady-state
tracking bound, and multi-fault feasibility.

\begin{theorem}[Hybrid practical-ISS with explicit recovery
  envelope under $F$ sequential agent failures]
\label{thm:hybrid-stability}
Under hypotheses H1--H3 (\cref{sec:stability-setup}), the
cable-tilt domain condition of \cref{lem:tension-cancel}, any
fault schedule $\sigma = \{(t_k^\star, i_k^\star)\}_{k=1}^{F}
\in \mathcal{F}$ of $F$ unannounced cable severances satisfying
$F \le N - 2$ and the per-fault actuator-margin condition
$m_L g/(N{-}k) \le \kappa_{\mathrm{act}} f_{\max}$ for all
$k \in \{1,\dots,F\}$, and admissible disturbances
$d \in \mathcal{D}$, the stacked tracking-and-tension-error
state $\vect\xi = (\vect e_p, \vect e_v, \vect\varepsilon_T)$
satisfies
\begin{align}
  \|\vect\xi(t)\|
    &\;\le\;
      \beta\bigl(\|\vect\xi(0)\|,\,t\bigr)
      + \gamma_w\bigl(\|w\|_\infty\bigr)
      + \gamma_q\bigl(\|q - q^\star\|_\infty\bigr)
      \notag \\
    &\quad
      + \gamma_\eta\bigl(\eta\bigr)
      + \sum_{k=1}^{F}
           \Gamma_f\,e^{-\lambda(t - t_k^\star)}\,
           \mathbf{1}_{t \ge t_k^\star},
  \label{eq:iss-envelope-explicit}
  \quad t \in [t_{\mathrm{pickup}},\, T],
\end{align}
where $\beta \in \mathcal{KL}$, $\gamma_w, \gamma_q,
\gamma_\eta \in \mathcal{K}$, $\lambda \triangleq
\alpha_{\min}/2 = \alpha_{xy}/2 = \SI{1.19}{\second^{-1}}$
is the norm-level recovery rate induced by the slower of the
altitude and horizontal closed-loop spectral rates
(\cref{lem:altitude-iss,lem:horizontal-iss}), $\eta =
\eta(\vect\xi[\cdot])$ is the realized slack duty cycle bounded
by $\eta_{\max}$ via H1, and $\Gamma_f > 0$ is the
fault-overshoot coefficient set by the redistribution geometry.
Equivalently, $V(\vect\xi(t_{k+1}^{\star,-})) \le
\rho\,V(\vect\xi(t_k^{\star,-})) + c$ for explicit $c > 0$
per \cref{lem:dwell-decay}. Each surviving cable's tension
obeys the practical envelope
\begin{equation}
  |T_i(t) - T_i^{\mathrm{qs}}(t)|
    \;\le\;
    T_{\mathrm{overshoot}}\,
    e^{-\lambda(t - t^\star_{\mathrm{last}})}
    + T_{\mathrm{floor}},
  \qquad i \in \mathcal{S}(t),
  \label{eq:tension-envelope-thm}
\end{equation}
where $T_{\mathrm{overshoot}}$ is bounded by the corresponding
tension-error radius extracted from \eqref{eq:chi-explicit},
and $T_{\mathrm{floor}}$ is the steady-state floor induced by
$\gamma_w$, $\gamma_q$, and $\gamma_\eta$.
\end{theorem}

\begin{proof}
We organize the argument in three parts: inter-fault ISS,
per-fault jump, and inductive composition.

\paragraph{Part 1: ISS between fault events.}
Fix an inter-fault interval $[t_k^{\star,+}, t_{k+1}^{\star,-}]$
during which $\mathcal{S}(t)$ is constant with
$N_s = N - k$ surviving cables. By \cref{lem:tension-cancel},
the altitude dynamics reduce to \eqref{eq:alt-error-ode}
with bounded residual $r_i(t)$ on $\Omega_\tau^{\mathrm{dwell}}$.
\Cref{lem:altitude-iss,lem:horizontal-iss} then give, for
the composite Lyapunov function
$V = \lambda_{\min}(P_v)^{-1}V_z + \lambda_{\min}(P_{xy})^{-1}V_{xy}
+ \kappa_V V_{\mathrm{tension}}$:
\begin{equation}
  \dot V
    \;\le\; -\alpha_{\min}\,V
    \;+\; \gamma_w(\|w\|_\infty)
    \;+\; \gamma_q(\|q - q^\star\|_\infty)
    \;+\; \gamma_\eta(\eta),
  \label{eq:combined-iss}
\end{equation}
with $\gamma_\eta$ absorbing the reduction residual
$C_1\delta + C_2\eta_{\max}$ of \cref{thm:reduction} and
$\alpha_{\min} = \alpha_{xy}$ from \cref{lem:horizontal-iss}.
By the standard comparison lemma,
integrating \eqref{eq:combined-iss} over any interval
$[0, t]$ yields
\begin{equation}
  V(t) \;\le\;
    e^{-\alpha_{\min} t} V(0)
    + \frac{1}{\alpha_{\min}}\bigl[
      \gamma_w(\|w\|_\infty)
      + \gamma_q(\|q - q^\star\|_\infty)
      + \gamma_\eta(\eta)\bigr],
\end{equation}
from which $\beta \in \mathcal{KL}$ and
$\gamma_w, \gamma_q, \gamma_\eta \in \mathcal{K}$ follow by
the standard ISS reformulation. Since $V$ is quadratic in
$\xi$ (with weighting matrices $P_v$, $P_{xy}$, $\kappa_V I$),
the Lyapunov-level decay at rate $\alpha_{\min}$ implies the
norm-level rate $\lambda = \alpha_{\min}/2$ via $V \sim
\|\xi\|^2$ (see \cref{rem:jump-vs-rate}).

\paragraph{Part 2: Lyapunov jump at each $t_k^\star$.}
By \cref{lem:jump-continuity}, $V(t_k^{\star,+}) \le
V(t_k^{\star,-}) + \chi(\Delta_{f,k})$ with $\chi$ bounded by
\eqref{eq:chi-explicit}. This bounded jump is then damped by
Part~1: integrating \eqref{eq:combined-iss} from
$t_k^{\star,+}$ over a dwell interval of length
$\tau_{d,k} \ge \tau_{\mathrm{pend}}$ gives
$V(t_{k+1}^{\star,-}) \le \rho\,V(t_k^{\star,-}) + c$ per
\cref{lem:dwell-decay}, with $\rho$ and $c$ as in
\eqref{eq:rho-def}--\eqref{eq:c-def}. The overshoot
$\chi(\Delta_{f,k})$ contributes to $c$ and decays
exponentially at the Lyapunov rate $\alpha_{\min}$; after
conversion through quadratic equivalence this produces
the term $\Gamma_f\,e^{-\lambda(t-t_k^\star)}$ in
\eqref{eq:iss-envelope-explicit}.

\paragraph{Part 3: Induction over $F$ fault events.}
Let $\rho < 1$ (from \cref{lem:dwell-decay}) and
suppose \eqref{eq:iss-envelope-explicit} holds for the first
$k-1$ faults. After the $k$-th fault,
$V(t_k^{\star,-}) \le \rho^{k-1}V(t_1^{\star,-}) + c\,(1 -
\rho^{k-1})/(1-\rho)$ by geometric summation, which remains
bounded for all $k \le F \le N-2$ as long as the
actuator-margin condition $m_L g/(N-k) \le \kappa_{\mathrm{act}}
f_{\max}$ is satisfied (ensuring $T_i^{\mathrm{qs}}(t_k^{\star,+})$
remains within the thrust envelope and $\chi(\Delta_{f,k})$
is finite). Composing over $F$ events via the
Teel--Goebel--Sanfelice hybrid-system framework~%
\cite[Theorem~1]{GoebelSanfeliceTeel2012} yields
\eqref{eq:iss-envelope-explicit}. The tension envelope
\eqref{eq:tension-envelope-thm} follows from
the positive-definite quadratic relation between
$V_{\mathrm{tension}}$ and $\|\vect\varepsilon_T\|^2$, together
with the exponential bound on $V_{\mathrm{tension}}$ from
Part~1, and the identity
$\varepsilon_{T,i} = T_i - T_i^{\mathrm{qs}}$; taking the
required square root yields the norm-level rate $\lambda$.
\end{proof}

\begin{figure}[t]
\centering
\begin{tikzpicture}[
    >={Stealth[length=2mm,width=1.6mm]},
    axislab/.style={font=\footnotesize},
    tick/.style={font=\scriptsize},
    fault/.style={dashed, red!70!black, thick},
    flow/.style={thick, blue!60!black},
    jump/.style={thick, red!70!black},
    annot/.style={font=\scriptsize}
  ]

  % Plot occupies x in [0, 8] s.  Top panel translated up.

  % ===== TOP PANEL: survivor cardinality |S(t)| ==============
  \begin{scope}[yshift=2.6cm]
    \draw[->, thick] (-0.2, 0) -- (8.4, 0)
          node[axislab, below right=-1pt] {$t$};
    \draw[->, thick] (0, -0.2) -- (0, 1.85)
          node[axislab, above right=-1pt] {$|\mathcal{S}(t)|$};

    \draw[flow] (0.0, 1.5) -- (2.0, 1.5);
    \draw[flow, dotted] (2.0, 1.5) -- (2.0, 1.0);
    \draw[flow] (2.0, 1.0) -- (5.0, 1.0);
    \draw[flow, dotted] (5.0, 1.0) -- (5.0, 0.5);
    \draw[flow] (5.0, 0.5) -- (8.0, 0.5);

    \node[tick, anchor=east] at (+1.0, 1.7) {$N{=}5$};
    \node[tick, anchor=east] at (+3.0, 1.2) {$N{=}4$};
    \node[tick, anchor=east] at (+6.0, 0.7) {$N{=}3$};

    \draw[fault] (2.0, -0.1) -- (2.0, 1.85);
    \draw[fault] (5.0, -0.1) -- (5.0, 1.85);
    \node[tick, red!70!black, anchor=south] at (2.0, 1.85)
          {$t_1^\star$};
    \node[tick, red!70!black, anchor=south] at (5.0, 1.85)
          {$t_2^\star$};
  \end{scope}

  % ===== BOTTOM PANEL: Lyapunov sawtooth V(t) ================
  \draw[->, thick] (-0.2, 0) -- (8.4, 0)
        node[axislab, below right=-1pt] {$t$};
  \draw[->, thick] (0, -0.2) -- (0, 2.45)
        node[axislab, above right=-1pt]
        {$V(\boldsymbol{\xi}(t))$};

  % flow segment 1
  \draw[flow]
    plot[smooth, samples=40, domain=0:2]
      (\x, {0.4 + 0.9*exp(-0.7*\x)});
  % jump 1
  \draw[jump, ->] (2.0, {0.4 + 0.9*exp(-1.4)})
                -- (2.0, 1.55);
  % flow segment 2
  \draw[flow]
    plot[smooth, samples=40, domain=2:5]
      (\x, {0.4 + 1.15*exp(-0.7*(\x-2))});
  % jump 2
  \draw[jump, ->] (5.0, {0.4 + 1.15*exp(-2.1)})
                -- (5.0, 1.10);
  % flow segment 3
  \draw[flow]
    plot[smooth, samples=40, domain=5:8]
      (\x, {0.4 + 0.7*exp(-0.7*(\x-5))});

  % faults — bottom labels
  \draw[fault] (2.0, -0.1) -- (2.0, 2.45);
  \draw[fault] (5.0, -0.1) -- (5.0, 2.45);
  \node[tick, red!70!black, anchor=north] at (2.0, -0.1)
        {$t_1^\star$};
  \node[tick, red!70!black, anchor=north] at (5.0, -0.1)
        {$t_2^\star$};

  % H2 dwell brace below the bottom panel
  \draw[decorate, decoration={brace, amplitude=2pt},
        thin] (2.0, 4.5) -- (5.0, 4.5)
        node[annot, midway, above=2pt]
        {$\tau_{d,1} \ge \tau_{\mathrm{pend}}$ (H2)};

  % ===== Short on-figure labels =============================
  % Each label is a single word + lemma tag.  Math expressions
  % are in the caption.  Labels live in the upper band, well
  % above the curves (the curves stay below y=1.6).

  % FLOW label — placed above the first decaying segment.
  \node[annot, blue!60!black, anchor=south]
        at (1.0, 1.95)
        {\textbf{flow} (Lem.~\ref{lem:altitude-iss})};

  % JUMP label — placed slightly to the right of jump 1 in
  % white space above the second decaying segment.
  \node[annot, red!70!black, anchor=south]
        at (3.5, 1.85)
        {\textbf{jump} (Lem.~\ref{lem:jump-continuity})};

  % CONTRACTION label — above the third flow segment.
  \node[annot, red!70!black, anchor=south]
        at (6.5, 1.55)
        {\textbf{contraction} (Lem.~\ref{lem:dwell-decay})};

  % chi labels — small white-backed labels left of each jump
  \node[annot, red!70!black, anchor=east, fill=white,
        inner sep=1.2pt] at (1.95, 1.10)
        {$\chi(\Delta_{f,1})$};
  \node[annot, red!70!black, anchor=east, fill=white,
        inner sep=1.2pt] at (4.95, 0.80)
        {$\chi(\Delta_{f,2})$};

\end{tikzpicture}
\caption{Hybrid sawtooth structure of the closed-loop Lyapunov
function across two unannounced cable-severance events.
\emph{Top:} the survivor set $\mathcal{S}(t)$ steps down by one
at each fault $t_k^\star$; cardinality is right-continuous.
\emph{Bottom:} $V(\boldsymbol{\xi})$ exhibits three regimes,
labelled inline. \textbf{Flow:} between events the continuous-time
decay of \cref{lem:altitude-iss} gives
$\dot V \le -\lambda V + \gamma_w(\|w\|) +
\gamma_\theta(\|\theta-\theta^\ast\|) + \gamma_\eta(\eta)$.
\textbf{Jump:} at each fault $t_k^\star$ the post-redistribution
geometry term satisfies
$V(t_k^{\star,+}) \le V(t_k^{\star,-}) + \chi(\Delta_f)$
(\cref{lem:jump-continuity}). \textbf{Contraction:} composing
flow with jumps over one inter-fault dwell yields the
per-cycle bound
$V(t_{k+1}^-) \le \rho\,V(t_k^-) + c$ with
$\rho = \exp(-\alpha_{\min}\tau_{\mathrm{pend}}) \approx 0.005$
(\cref{lem:dwell-decay}). The H2 hypothesis
$\tau_{d,k} \ge \tau_{\mathrm{pend}}$ is structural: it is the
minimum window over which the feed-forward identity
$T_i^{\mathrm{ff}} = T_i$ contracts the post-jump overshoot
back below $V(t_k^-)$.}
\label{fig:lyapunov-sawtooth}
\end{figure}


\begin{remark}[Practical-stability interpretation]
\Cref{thm:hybrid-stability} is a practical stability bound,
not asymptotic convergence. The ISS framework treats the
reference waypoints as a time-varying input, not a
disturbance; the residual captured by the $\gamma$ terms
reflects the reference's dynamic content---principally the
quasi-static pendulum offset---rather than a closed-loop
failure. Within the analyzed baseline architecture and gain schedule,
$T_i^{\mathrm{ff}} = T_i$ supplies the bounded-jump mechanism
that makes \eqref{eq:iss-envelope-explicit} applicable; no
detection or inter-drone communication is required.
\end{remark}

\begin{remark}[Jump-bound vs.\ contraction-rate mechanism]
\label{rem:jump-vs-rate}
Two distinct mechanisms produce the dwell-cycle contraction.
The \emph{jump-bound mechanism} is the identity
$T_i^{\mathrm{ff}} = T_i$, whose contribution is to reduce the
size of the bounded reset observed in
\cref{lem:jump-continuity}; the continuity of the tracking
coordinates follows from state continuity at the severance.
The
\emph{contraction-rate mechanism} is the closed-loop
spectrum encoded in $\alpha_{xy}$ of
\cref{lem:horizontal-iss}; the inter-event flow contracts
$V$ at this rate. The ablation of \cref{sec:VI-D} supports
the first mechanism empirically: disabling the identity
increases peak sag by $3.6$--$4.0\times$ and produces larger
post-fault Lyapunov jumps in \cref{fig:fault-zoom}. The
jump analysis uses only the bounded-jump property; the exact
inflation without feed-forward is scenario-dependent. The
architectural innovation is the jump-bound mechanism, while
$\alpha_{xy}$ determines the Lyapunov contraction rate and,
through quadratic equivalence, the induced norm recovery rate
$\lambda = \alpha_{xy}/2$ that governs how quickly each
bounded jump is forgotten between faults.
\end{remark}

% ---- Corollaries (unchanged except minor notation) --------
\begin{corollary}[Pre-fault steady-state tracking bound]
\label{cor:steady-state}
Prior to any fault event, the payload tracking error is
dominated by the quasi-static pendulum offset induced by the
reference acceleration, and admits the explicit envelope
\begin{equation}
  \bar{e}_{p,xy} \le \frac{L A_{xy,\max}}{g}
    \left(1 - \frac{1}{2\zeta^2}\right)
    + \frac{A_{xy,\max}}{\kappa_{\mathrm{QP}} K_p^{xy}}
    + O(\delta),
  \label{eq:ss-bound}
\end{equation}
where $A_{xy,\max}$ is the maximum reference acceleration,
$\zeta = 1.12$ from \cref{lem:antiswing-invariant}, and
$\kappa_{\mathrm{QP}} = w_t/(w_t + w_e) \approx 0.98$ the
QP regularization factor. At canonical $A_{xy,\max} \approx
\SI{2.47}{\metre\per\second\squared}$,
\begin{align}
  \bar{e}_{p,xy}
    \;\lesssim\;
    \frac{1.25 \times 2.47}{9.81}\!
    &\left(1 - \frac{1}{2\times 1.12^2}\right)
      + \frac{2.47}{0.98 \times 30}
    \approx 0.19 + 0.08
    \approx \SI{0.27}{\metre}.
  \label{eq:ss-numeric}
\end{align}
\eqref{eq:ss-numeric} is a leading-order scale, not a strict
envelope: the first term uses a low-frequency approximation
of the pendulum transfer-function bias.
\end{corollary}

\begin{proof}[Derivation]
The pendulum transfer function settles to
$\bar{e}_{\mathrm{pend}} = (LA_{xy,\max}/g)(1 -
1/(2\zeta^2))$ at quasi-static frequency; the outer-loop PD
contributes $A_{xy,\max}/(\kappa_{\mathrm{QP}} K_p^{xy})$;
the reduction error is $O(\delta)$ by \cref{thm:reduction}.
Summing gives \eqref{eq:ss-bound}. V1 reports horizontal
RMSE $\SI{0.306}{\metre}$ vs.\ the analytical
$\SI{0.27}{\metre}$; the gap is the centripetal contribution
on the curved segments, which the quasi-static derivation
excludes.
\end{proof}

\begin{corollary}[Robustness to mass, stiffness, drag, and wind]
\label{cor:robustness}
The envelope \eqref{eq:iss-envelope-explicit} is robust along
four axes. \emph{Mass mismatch}: $T_i^{\mathrm{ff}} = T_i$
uses measured tension, making \cref{lem:tension-cancel,%
lem:altitude-iss} mass-independent to leading order.
\emph{Stiffness and drag mismatches}: enter through
$O(\delta + \eta_{\max})$, absorbed by $\gamma_\eta$ without
degrading $\rho$ or $\alpha_{xy}$. \emph{Wind}: enters
through $\gamma_w$ with steady-state gain $1/K_p^{xy}$.
The P2-B mass sweep yields only $1.9\%$ RMSE variation over
a $56\%$ mass range with FF active; V2 yields
$\SI{0.312}{\metre}$ RMSE under Dryden wind, matching
the no-wind V1.
\end{corollary}

\paragraph{Multi-fault feasibility.}
\eqref{eq:iss-envelope-explicit} extends to $k$ sequential
faults by induction under two conditions: (a)~geometric
closure $k \le N - 2$, (b)~actuator margin
$m_L g/(N{-}k) \le \kappa_{\mathrm{act}} f_{\max}$ at every
step. Canonical $N = 5$, $m_L = \SI{10}{\kilo\gram}$ admits
$k = 2$ with margin; V4 and V5 (dual-fault) yield
3-D RMSE $0.324$--$\SI{0.328}{\metre}$, confirming
\cref{thm:hybrid-stability}.

\paragraph{Empirical validation summary.}
\label{sec:stability-empirical-validation}
\Cref{tab:stability-evidence} maps each claim to its campaign
observation; full context is in \cref{sec:simulation}.

\begin{table*}[t]
  \centering
  \caption{Empirical support for hybrid practical-ISS claims.}
  \label{tab:stability-evidence}
  \small
  \renewcommand{\arraystretch}{1.05}
  \begin{tabular}{p{5.5cm}p{2.5cm}p{8cm}}
    \toprule
    Claim / Lemma & Campaign & Evidence \\
    \midrule
    Tension cancellation (\cref{lem:tension-cancel})
      & P2-A ablation
      & $3.6\times$--$4.0\times$ sag increase with FF
        disabled \\
    Altitude ISS decay (\cref{lem:altitude-iss})
      & V1 nominal
      & 3-D RMSE $= \SI{0.312}{\metre}$, stable;
        $\lambda_z = \SI{5.37}{\second^{-1}}$ from $P_v$
        in \eqref{eq:Pv-numeric} \\
    Horizontal ISS decay (\cref{lem:horizontal-iss})
      & V1--V5 series
      & $\alpha_{xy} = \SI{2.38}{\second^{-1}}$
        from $P_{xy}$ in \eqref{eq:Pxy-numeric};
        $\rho = \exp(-\alpha_{xy}\tau_{\mathrm{pend}})
        \approx 0.005$ \\
    Anti-swing invariance (\cref{lem:antiswing-invariant})
      & V1--V5 series
      & Post-fault RMSE within $4\%$ of nominal
        across all variants \\
    Bounded fault jump (\cref{lem:jump-continuity})
      & V3, V4, V5 traces
      & $\hat\chi$ peak $1.1\!\times\!10^{-4}$ (V4),
        $1.3\!\times\!10^{-5}$ (V5); qualitatively consistent
        with the bounded redistribution estimate
        \eqref{eq:chi-explicit}; under equal-sharing hover
        the canonical coefficient reduces to
        $\kappa_V\Delta_f^2/(2N_s)$, with quantitative
        calibration of $\kappa_V$ deferred to follow-up work \\
    Dwell-cycle contraction $\rho < 1$
        (\cref{lem:dwell-decay})
      & V4, V5 dual-fault
      & $\hat\rho_{\mathrm{V4}} = 0.20$,
        $\hat\rho_{\mathrm{V5}} = 0.045$;
        both satisfy \eqref{eq:rho-hat-decomp}
        with $c/V(t_k^{\star,-})$ set by the Dryden
        disturbance floor (\cref{rem:rho-gap}) \\
    Mass robustness (\cref{cor:robustness})
      & P2-B sweep
      & $1.9\%$ RMSE range over $56\%$ mass variation \\
    Wind robustness (\cref{cor:robustness})
      & V2 wind
      & V2 3-D RMSE $= \SI{0.312}{\metre} \approx$ V1 \\
    \bottomrule
  \end{tabular}
\end{table*}

The baseline theorem delivers \eqref{eq:iss-envelope-explicit}
under H1--H3 and the identity $T_i^{\mathrm{ff}} = T_i$
within the analyzed architecture; no additional architecture
is required for \cref{def:problem}. \Cref{sec:extensions} develops parameter
adaptation, tension-ceiling enforcement, and post-fault
reshape---each targeting a residual channel of
\eqref{eq:iss-envelope-explicit}---and shows each preserves
H1--H3 on the baseline domain.